\chapter{Praxisprojekte -- Steigender Schwierigkeitsgrad}

\label{chap:praxisprojekte}

% ============================================================
% LERNZIELE
% ============================================================
\begin{lernziele}
Nach diesem Kapitel hast du:
\begin{itemize}
    \item Fünf vollständige KI-Projekte mit steigendem Schwierigkeitsgrad gebaut
    \item Einen OpenAI-Chatbot mit Streaming und Chat-History implementiert
    \item Ein RAG-System mit ChromaDB und semantischer Suche aufgebaut
    \item Ein Modell per Fine-Tuning auf eigene Daten angepasst
    \item Ein Multi-Agent-System mit LangGraph orchestriert
    \item Eine vollständige Produktions-Pipeline mit Monitoring deployt
\end{itemize}
\end{lernziele}

% ============================================================
% MOTIVATION
% ============================================================
\section{Motivation}

Learning by Doing. Theorie allein reicht nicht -- erst Praxis macht aus Wissen Können. Jedes Projekt baut auf dem vorherigen auf: Chatbot (Projekt 1), RAG (Projekt 2), Fine-Tuning (Projekt 3), Multi-Agent (Projekt 4), Produktions-Pipeline (Projekt 5). Projekt 6 als optionale Herausforderung.

\section{Grundlagen -- Projektstruktur und Setup}

Jedes Projekt folgt der gleichen Struktur:

\begin{itemize}[leftmargin=*]
    \item \textbf{Ziel}: Was du lernst und baust
    \item \textbf{Technische Anforderungen}: Benötigte Tools, Libraries, API-Keys
    \item \textbf{Code-Rahmen}: Startcode, den du erweitern kannst
    \item \textbf{Erweiterungsideen}: Wie du das Projekt professionell ausbaust
    \item \textbf{Typische Fallstricke}: Was schiefgehen kann und wie du es vermeidest
\end{itemize}

\textbf{Grundlegende Setup-Empfehlungen:}
\begin{itemize}[leftmargin=*]
    \item Python 3.12+ mit \texttt{uv} oder \texttt{pip}
    \item Virtuelle Umgebung pro Projekt (\texttt{uv venv} oder \texttt{python -m venv})
    \item API-Keys in \texttt{.env} Datei, nicht im Code
    \item Git-Repository von Tag 1
\end{itemize}

% ============================================================
% PROJEKT 1
% ============================================================
\section{Projekt 1: OpenAI Chatbot mit Streamlit}

\textbf{Schwierigkeit: }$\star$ $\star$ $\star$ ---
\textbf{Dauer: }~2--3 Stunden

\subsection*{Ziel}
Ein funktionsfähiger Chatbot, der die OpenAI Chat Completions API nutzt -- mit Echtzeit-Streaming und Chat-History.

\subsection*{Technische Anforderungen}
\begin{itemize}[leftmargin=*]
    \item Python 3.10+, openai-Paket, streamlit
    \item Streamlit für die UI (Chat-Interface)
    \item OpenAI API Key (\$5 Guthaben reicht für Testing)
    \item Chat-History mit Kontext-Limit-Monitoring
\end{itemize}

\subsection*{Code-Rahmen}

\begin{lstlisting}[language=Python, caption=Minimaler Streamlit-Chatbot -- das komplette Projekt in ~80 Zeilen]
import streamlit as st
import openai
from tiktoken import encoding_for_model

st.set_page_config(page_title="AI Chatbot", page_icon=":robot_face:")
st.title("Mein AI-Assistent")

if "messages" not in st.session_state:
    st.session_state.messages = []
if "api_key_set" not in st.session_state:
    st.session_state.api_key_set = False

if not st.session_state.api_key_set:
    api_key = st.text_input("OpenAI API Key:", type="password")
    if api_key:
        client = openai.OpenAI(api_key=api_key)
        st.session_state.api_key_set = True

if st.session_state.api_key_set:
    for msg in st.session_state.messages:
        with st.chat_message(msg["role"]):
            st.markdown(msg["content"])

    if prompt := st.chat_input("Schreibe eine Nachricht..."):
        st.session_state.messages.append({"role": "user", "content": prompt})

        encoding = encoding_for_model("gpt-4o-mini")
        total_tokens = sum(
            len(encoding.encode(m["content"]))
            for m in st.session_state.messages
        )

        if total_tokens > 120_000:
            st.warning(f"Kontext-Limit nahe! ({total_tokens}/128000 Tokens)")
            st.session_state.messages = st.session_state.messages[-5:]

        with st.chat_message("user"):
            st.markdown(prompt)

        with st.chat_message("assistant"):
            message_placeholder = st.empty()
            full_response = ""

            for chunk in client.chat.completions.create(
                model="gpt-4o-mini",
                messages=st.session_state.messages,
                stream=True,
                max_tokens=1024,
            ):
                if chunk.choices[0].delta.content:
                    full_response += chunk.choices[0].delta.content
                    message_placeholder.markdown(full_response + "...")

            message_placeholder.markdown(full_response)
            st.session_state.messages.append(
                {"role": "assistant", "content": full_response}
            )

st.sidebar.title("Info")
st.sidebar.info(f"Token-Verbrauch: {total_tokens:,} / 128000")
st.sidebar.text("Gewaehlt: gpt-4o-mini (EUR0,0015/1K Input)")
\end{lstlisting}

\subsection*{Erweiterungsideen}
\begin{itemize}[leftmargin=*]
    \item Modellauswahl in der Sidebar (gpt-4o-mini vs. gpt-4o vs. claude)
    \item System-Prompt-Konfiguration im Settings-Panel
    \item Temperature und max\_tokens als Slider
    \item Chat exportieren/importieren (JSON)
\end{itemize}

\subsection*{Typische Fallstricke}
\begin{itemize}[leftmargin=*]
    \item API-Key nicht gesetzt: Immer auf \texttt{st.session\_state} prüfen
    \item Token-Limit ignorieren: Ohne Kontext-Kompression steigen die Kosten linear
    \item Stream nicht anzeigen: Ohne streaming wirkt der Bot langsam
\end{itemize}

\begin{praxishinweis}
Streamlit eignet sich für Prototypen, nicht für Produktion. Für Produktivsysteme: FastAPI-Backend mit React/Vue-Frontend. Streamlit teilt einen Python-Prozess -- bei vielen Usern wirds eng.
\end{praxishinweis}

% ============================================================
% PROJEKT 2
% ============================================================
\section{Projekt 2: RAG-System über eigene Dokumente mit ChromaDB}

\textbf{Schwierigkeit: }$\star$ $\star$ $\star$ $\star$ ---
\textbf{Dauer: }~4--6 Stunden

\subsection*{Ziel}
System, das PDFs oder Textdokumente in einen Vector-Index umwandelt und semantisch abfragbar macht. Der Hello World des AI Engineers.

\subsection*{Technische Anforderungen}
\begin{itemize}[leftmargin=*]
    \item Python 3.10+, langchain, chromadb, openai
    \item PDF-Parsing (PyPDF2 oder pdfplumber)
    \item OpenAI Embeddings (text-embedding-3-small)
    \item Streamlit für die Query-Oberfläche
\end{itemize}

\subsection*{Code-Rahmen}

\begin{lstlisting}[language=Python, caption=RAG-System mit ChromaDB -- Dokumenten-Index und semantische Suche]
import streamlit as st
import os
from langchain.text_splitter import RecursiveCharacterTextSplitter
from langchain_chroma import Chroma
from langchain_openai import OpenAIEmbeddings, ChatOpenAI, ChatPromptTemplate
from langchain_core.output_parsers import StrOutputParser
from pdfplumber import open as open_pdf

st.set_page_config(page_title="RAG Document Q&A", layout="wide")
st.title("Dokumenten-Fragebeantwortung")

if "rag_chain" not in st.session_state:
    st.session_state.rag_chain = None
if "db" not in st.session_state:
    st.session_state.db = None

st.sidebar.header("Dokumente")
uploaded = st.sidebar.file_uploader(
    "PDF oder TXT hochladen", type=["pdf", "txt"]
)
if uploaded and st.sidebar.button("Index erstellen"):
    with st.spinner("Lese und indiziere..."):
        text = ""
        if uploaded.name.endswith(".pdf"):
            with open_pdf(uploaded) as pdf:
                text = "\n".join(
                    page.extract_text()
                    for page in pdf.pages
                    if page.extract_text()
                )
        else:
            text = uploaded.getvalue().decode("utf-8")

        splitter = RecursiveCharacterTextSplitter(
            chunk_size=512, chunk_overlap=64,
            separators=["\n\n", "\n", " ", ""]
        )
        chunks = splitter.split_text(text)

        embeddings = OpenAIEmbeddings(model="text-embedding-3-small")
        db = Chroma(
            collection_name="user_docs",
            embedding_function=embeddings,
            persist_directory="./chroma_db"
        )
        db.add_texts(chunks)
        st.session_state.db = db

        retriever = db.as_retriever(search_kwargs={"k": 5})

        prompt = ChatPromptTemplate.from_messages([
            ("system", "Beantworte die Frage basierend auf den folgenden "
                       "Dokumenten. Wenn du es nicht weisst: "
                       "'Dazu habe ich keine Daten.' Zitiere die Quelle."),
            ("human", "Relevante Abschnitte:\n{context}\n\nFrage: {question}")
        ])

        llm = ChatOpenAI(model="gpt-4o-mini", temperature=0.1)
        rag_chain = (
            {"context": retriever, "question": lambda x: x}
            | prompt | llm | StrOutputParser()
        )
        st.session_state.rag_chain = rag_chain

        st.sidebar.success("Index erstellt!")

if st.session_state.db and (
    prompt := st.text_input(
        "Stelle deine Frage:",
        placeholder="Was steht im Dokument ueber ...?"
    )
):
    with st.spinner("Durchsuche Dokumente..."):
        answer = st.session_state.rag_chain.invoke(prompt)

    st.markdown("---")
    st.subheader("Antwort:")
    st.markdown(answer)
\end{lstlisting}

\subsection*{Erweiterungsideen}
\begin{itemize}[leftmargin=*]
    \item Multiple Dokumente parallel indexieren
    \item Filter nach Dokument-Typ (PDF, TXT, DOCX)
    \item Hybride Suche mit BM25
    \item Export der Chat-History als PDF
\end{itemize}

\subsection*{Typische Fallstricke}
\begin{itemize}[leftmargin=*]
    \item Chunking-Grösse falsch gewählt: Zu kleine Chunks verlieren Kontext, zu grosse passen nicht ins Context Window
    \item Embedding-Modell wechseln: Ändert alle bestehenden Vektoren. Einmal gewählt, bleibt es
    \item Kein persistierender Speicher: ChromaDB in-memory verliert Daten beim Neustart
\end{itemize}

\begin{praxishinweis}
Die Chunking-Strategie entscheidet über RAG-Qualität. Bewährt: 256--512 Tokens pro Chunk mit 10--20\% Overlap. Semantische Chunks (Absätze) besser als reine Token-Zerlegung. Embedding-Modell nicht nach dem Indexieren wechseln.
\end{praxishinweis}

% ============================================================
% PROJEKT 3
% ============================================================
\section{Projekt 3: Fine-Tuning eines Modells auf Firmendaten}

\textbf{Schwierigkeit}: $\star$ $\star$ $\star$ $\star$ ---
\textbf{Dauer}: ~6--8 Stunden (inkl. Daten-Sammlung)

\subsection*{Ziel}
Ein GPT-4o-mini-Modell auf Firmendaten fine-tunen -- z.B. für einen spezifischen Antwort-Stil, Terminologie oder Format.

\subsection*{Technische Anforderungen}
\begin{itemize}[leftmargin=*]
    \item 50--200 qualitativ hochwertige Input-Output-Paare
    \item JSONL-Format für die Trainingsdaten
    \item OpenAI Fine-Tuning API
    \item Evaluation gegen ein Golden Dataset
\end{itemize}

\subsection*{Code-Rahmen}

\begin{lstlisting}[language=Python, caption=Fine-Tuning Pipeline -- Trainingsdaten erstellen und Modell trainieren]
import json
import openai

# === Schritt 1: Trainingsdaten erstellen (JSONL) ===
training_data = [
    {
        "messages": [
            {"role": "user", "content": "Was kostet das Enterprise-Ticket?"},
            {"role": "assistant", "content": "Das Enterprise-Ticket kostet "
             "99 EUR/Monat exkl. MwSt."}
        ]
    },
    {
        "messages": [
            {"role": "user", "content": "Kann ich das Abo kuendigen?"},
            {"role": "assistant", "content": "Ja, Sie koennen Ihr Abo "
             "jederzeit kuendigen. Gehen Sie zu Einstellungen > Abonnement."}
        ]
    },
]

with open("fine_tuning_data.jsonl", "w") as f:
    for item in training_data:
        f.write(json.dumps(item) + "\n")

# === Schritt 2: Datei hochladen ===
client = openai.OpenAI()
file_response = client.files.create(
    file=open("fine_tuning_data.jsonl", "rb"),
    purpose="fine-tune"
)

# === Schritt 3: Fine-Tuning Job starten ===
ft_job = client.fine_tuning.jobs.create(
    training_file=file_response.id,
    model="gpt-4o-mini",
    hyperparameters={"n_epochs": 3, "learning_rate_multiplier": 0.1},
)

print(f"Fine-Tuning gestartet! Job ID: {ft_job.id}")

# === Schritt 4: Evaluieren ===
def eval_fine_tuned(model_name, test_data):
    for example in test_data[:5]:
        baseline = client.chat.completions.create(
            model="gpt-4o-mini",
            messages=[{"role": "user",
                       "content": example["input"]}]
        ).choices[0].message.content

        fine_tuned = client.chat.completions.create(
            model=model_name,
            messages=[{"role": "user",
                       "content": example["input"]}]
        ).choices[0].message.content

        print(f"Input: {example['input'][:40]}...")
        print(f"Baseline:   {baseline[:60]}...")
        print(f"FineTuned:  {fine_tuned[:60]}...")
        print("---")
\end{lstlisting}

\subsection*{Erweiterungsideen}
\begin{itemize}[leftmargin=*]
    \item Automatische Datengenerierung mit GPT-4o für Daten-Augmentation
    \item Hyperparameter-Tuning via Grid Search
    \item A/B Testing: Fine-Tuned Model vs. Baseline gegen Golden Dataset
\end{itemize}

\subsection*{Typische Fallstricke}
\begin{itemize}[leftmargin=*]
    \item Zu wenig Daten: 50 Beispiele sind das absolute Minimum. 150+ für gute Ergebnisse
    \item Datenqualität: Garbage in, Garbage out. Schlechte Beispiele verschlechtern das Modell
    \item Overfitting: Zu viele Epochen (n\_epochs > 10) führen zu Overfitting
\end{itemize}

\begin{praxishinweis}
Fine-Tuning lohnt bei spezifischem Output-Format oder Terminologie. Für Wissens-Update ist RAG der bessere Weg. Faustregel: < 300 Beispiele? Prompt-Engineering + RAG. Mindestens 100 hochwertige Paare für Fine-Tuning.
\end{praxishinweis}

% ============================================================
% PROJEKT 4
% ============================================================
\section{Projekt 4: Multi-Agent System mit LangGraph}

\textbf{Schwierigkeit}: $\star$ $\star$ $\star$ $\star$ $\star$ ---
\textbf{Dauer}: ~12--16 Stunden

\subsection*{Ziel}
Ein System mit mehreren spezialisierten Agenten, die zusammenarbeiten -- ein Recherche-Agent, ein Analyse-Agent und ein Writer-Agent -- orchestriert via LangGraph.

\subsection*{Technische Anforderungen}
\begin{itemize}[leftmargin=*]
    \item Python 3.10+, langgraph, langchain, openai
    \item State-Management für den Graph Workflow
    \item Paralleles Tool Calling für die Agent-Kommunikation
    \item Error Handling und Fallback-Mechanismen
\end{itemize}

\subsection*{Code-Rahmen}

\begin{lstlisting}[language=Python, caption=Multi-Agent-System mit LangGraph]
from langgraph.graph import StateGraph, MessagesState, END
from typing import TypedDict

class AgentState(TypedDict):
    messages: list[dict]
    research_done: bool
    analysis_ready: bool

def researcher_node(state: AgentState) -> AgentState:
    """Agent 1: Recherchiert in der Wissensdatenbank."""
    last_msg = state["messages"][-1]["content"]
    research_results = search_knowledgebase(last_msg)
    return {
        "messages": [{"role": "assistant",
                      "content": f"[Forschungsergebnis]\n{research_results}"}],
        "research_done": True,
    }

def analyst_node(state: AgentState) -> AgentState:
    """Agent 2: Analysiert die Recherche-Ergebnisse."""
    research_text = "\n".join(
        m["content"] for m in state["messages"]
        if "[Forschungsergebnis]" in m.get("content", "")
    )
    analysis = client.chat.completions.create(
        model="gpt-4o",
        messages=[{"role": "user",
                   "content": f"Analysiere:\n\n{research_text}"}]
    )
    return {
        "messages": [{"role": "assistant",
                      "content": f"[Analyse]\n{analysis.choices[0].message.content}"}],
        "analysis_ready": True,
    }

def writer_node(state: AgentState) -> AgentState:
    """Agent 3: Schreibt den finalen Bericht."""
    all_content = "\n".join(m["content"] for m in state["messages"])
    final_report = client.chat.completions.create(
        model="gpt-4o",
        messages=[{"role": "user",
                   "content": f"Schreibe Bericht basierend auf:\n\n{all_content}"}]
    )
    return {
        "messages": [{"role": "assistant",
                      "content": final_report.choices[0].message.content}],
    }

workflow = StateGraph(AgentState)
workflow.add_node("researcher", researcher_node)
workflow.add_node("analyst", analyst_node)
workflow.add_node("writer", writer_node)
workflow.set_entry_point("researcher")
workflow.add_edge("researcher", "analyst")
workflow.add_edge("analyst", "writer")
workflow.add_edge("writer", END)

app = workflow.compile()
result = app.invoke({
    "messages": [{"role": "user",
                  "content": "Analysiere die neuesten KI-Trends."}],
})
\end{lstlisting}

\subsection*{Erweiterungsideen}
\begin{itemize}[leftmargin=*]
    \item Bedingte Kanten: Wenn Recherche unzureichend, zurück zum Researcher
    \item Human-in-the-Loop bei kritischen Entscheidungen
    \item Persistenter State in Redis für langlaufende Workflows
\end{itemize}

\begin{praxishinweis}
Multi-Agent-Systeme multiplizieren Kosten und Fehlerquellen. Jeder Agent ruft mindestens einmal das LLM auf. Ohne Budget-Limits und Step-Count-Begrenzung laufen Kosten davon. Starte mit max. 3 Agenten und harter Step-Obergrenze.
\end{praxishinweis}

% ============================================================
% PROJEKT 5
% ============================================================
\section{Projekt 5: Evaluations-Dashboard mit DeepEval}

\textbf{Schwierigkeit}: $\star$ $\star$ $\star$ $\star$ ---
\textbf{Dauer}: ~6--8 Stunden

\subsection*{Ziel}
Ein automatisiertes Evaluations-System, das einen Chatbot gegen ein Golden Dataset testet, Metriken sammelt und Reports generiert. Die Grundlage für datengetriebene Prompt-Entwicklung.

\subsection*{Technische Anforderungen}
\begin{itemize}[leftmargin=*]
    \item Python 3.10+, deepeval, pytest, langfuse
    \item Golden Dataset mit 50+ Testfällen
    \item LLM-as-a-Judge (GPT-4o oder Claude)
    \item Report-Generierung (HTML/Markdown)
\end{itemize}

\subsection*{Code-Rahmen}

\begin{lstlisting}[language=Python, caption=Evaluations-Dashboard mit DeepEval]
from deepeval import evaluate
from deepeval.test_case import LLMTestCase
from deepeval.metrics import (
    AnswerRelevancyMetric,
    FaithfulnessMetric,
    HallucinationMetric,
    GEval
)
from deepeval.dataset import GoldenDataset

# Golden Dataset definieren
dataset = GoldenDataset(
    name="customer-support",
    test_cases=[
        LLMTestCase(
            input="Wie lautet eure Rueckgabepolitik?",
            actual_output="Sie haben 30 Tage Rueckgaberecht.",
            expected_output="30 Tage Rueckgaberecht, kostenlos",
            context=["Rueckgabepolitik: 30 Tage, kostenlos"]
        ),
    ]
)

# Metriken definieren
metrics = [
    AnswerRelevancyMetric(threshold=0.7),
    FaithfulnessMetric(threshold=0.8),
    HallucinationMetric(threshold=0.3),
    GEval(
        name="Professionalism",
        criteria="Bewerte die Professionalitaet der Antwort",
        evaluation_steps=[
            "Ist die Antwort hoeflich?",
            "Ist die Antwort vollstaendig?",
            "Enthaelt sie Quellenangaben?"
        ]
    )
]

# Evaluierung laufen lassen
results = evaluate(
    test_cases=dataset.test_cases,
    metrics=metrics,
    run_async=False
)

# Report generieren
for test_case in results.test_cases:
    print(f"\nInput: {test_case.input[:40]}...")
    for metric in test_case.metrics:
        status = "BESTANDEN" if metric.success else "FAILED"
        print(f"  [{status}] {metric.__name__}: "
              f"{metric.score:.2f}")
\end{lstlisting}

\subsection*{Erweiterungsideen}
\begin{itemize}[leftmargin=*]
    \item CI/CD-Integration: GitHub Action, die bei jedem Push evaluiert
    \item Dashboard: Streamlit-App, die Evaluations-Historie zeigt
    \item Regression-Alarme: Slack-Benachrichtigung bei Score-Abfall
\end{itemize}

\begin{praxishinweis}
Ein Golden Dataset ist das wichtigste Eval-Werkzeug. Starte mit 20--30 Fällen, erweitere auf 100+. Jeder Prompt-Change läuft gegen das Dataset. LLM-as-a-Judge ist günstig, aber biased (bevorzugt lange, eigene Antworten). Kreuzvalidierung mit verschiedenen Judge-Modellen empfohlen.
\end{praxishinweis}

% ============================================================
% PROJEKT 6
% ============================================================
\section{Projekt 6: Vollständige Produktions-Pipeline mit Monitoring}

\textbf{Schwierigkeit}: $\star$ $\star$ $\star$ $\star$ $\star$ ---
\textbf{Dauer}: ~20--30 Stunden

\subsection*{Ziel}
Eine vollständige, produktionsreife LLM-Anwendung mit allen Konzepten aus diesem Buch -- Chatbot + RAG + Agent + Evaluation + Monitoring. Das ultimative Portfolio-Projekt.

\subsection*{Architektur-Stack}
\begin{itemize}[leftmargin=*]
    \item FastAPI als Backend (Chat-Endpoint + Health-Check)
    \item ChromaDB persistent für RAG-Daten
    \item Redis für Prompt-Caching und Session-State
    \item Langfuse für Tracing und Observability
    \item Docker für Containerisierung
    \item Streamlit oder React für die Frontend-Oberfläche
    \item Prometheus + Grafana für Metriken
\end{itemize}

\subsection*{Architektur-Diagramm}

\begin{verbatim}
+------------------+       +------------------+
|   Streamlit UI   |------>|   FastAPI        |
|   (Frontend)     |<------|   (Backend)      |
+------------------+       +--------+---------+
                                    |
                    +---------------+---------------+
                    |               |               |
              +-----v-----+  +-----v-----+  +------v------+
              |  Redis    |  |  ChromaDB |  |  Langfuse   |
              |  (Cache)  |  |  (RAG)    |  |  (Tracing)  |
              +-----------+  +-----------+  +-------------+
\end{verbatim}

\subsection*{Checkliste für ein produktionsfähiges Projekt}

\begin{enumerate}[leftmargin=*]
    \item [\ ] Prompt-Templates versioniert und in Config-Dateien gespeichert
    \item [\ ] Golden Dataset mit 20+ Testbeispielen vorhanden
    \item [\ ] Token-Counting vor jedem API-Call
    \item [\ ] max\_tokens für alle LLM-Aufrufe gesetzt
    \item [\ ] Streaming-Response implementiert
    \item [\ ] Rate Limiting und Fallback-Mechanismen eingebaut
    \item [\ ] Input Sanitization (PII Masking) implementiert
    \item [\ ] Error Handling mit Graceful Degradation
    \item [\ ] Langfuse Tracing für alle LLM-Aufrufe
    \item [\ ] Health-Endpoint für Load Balancer
    \item [\ ] Dockerfile und docker-compose.yml vorhanden
    \item [\ ] Prometheus-Metriken für Requests, Latenz, Kosten
    \item [\ ] README mit Architektur-Diagramm und Setup-Anleitung
    \item [\ ] Multi-Provider-Fallback (OpenAI + Anthropic)
    \item [\ ] Circuit-Breaker für API-Ausfälle
\end{enumerate}

\subsection*{Erweiterungsideen}
\begin{itemize}[leftmargin=*]
    \item A/B-Testing: Zwei Prompt-Varianten parallel ausliefern
    \item Canary-Deployment: Neue Version nur für 5\% der User
    \item Automatische Kosten-Alarme bei Überschreitung von Budgets
\end{itemize}

\begin{praxishinweis}
Beginne nicht mit Projekt 6. Es kombiniert alle Konzepte -- überwältigend. Baue erst Projekte 1-5 einzeln, dann integriere. Das Docker-Setup und CI/CD aus Projekt 6 lassen sich auf jedes Einzelprojekt anwenden.
\end{praxishinweis}

% ============================================================
% ZUSAMMENFASSUNG
% ============================================================
\begin{zusammenfassung}
\begin{itemize}[leftmargin=*]
    \item Die Projekte sind steigend komplex: Vom Chatbot bis zur Multi-Agent-Produktions-Pipeline
    \item Jeder Schritt baut auf dem vorherigen auf -- Projekt 1 wird zur Basis des Gesamtprojekts
    \item Die Checkliste in Projekt 6 ist dein Blueprint für ein portfolio-würdiges AI Engineer-Projekt
    \item Git + README + Docker = professionelles Portfolio
    \item Wähle ein Projekt und baue es -- nicht nur lesen, sondern machen
\end{itemize}
\end{zusammenfassung}

\begin{merke}
Der Unterschied zwischen einem Hobby-Projekt und einem Portfolio-Projekt liegt nicht in der Komplexität, sondern in der Professionalität: sauberer Code, Tests, Docker und README machen den Unterschied.
\end{merke}

\section{Praxisprojekt -- Eigenes Projekt aufbauen}

\textbf{Ziel:} Wähle eines der sechs Projekte aus, implementiere es vollständig und dokumentiere es auf GitHub.

\textbf{Aufgaben:}
\begin{enumerate}[leftmargin=*]
    \item Wähle ein Projekt basierend auf deinem aktuellen Kenntnisstand
    \item Implementiere den Code-Rahmen und erweitere ihn
    \item Schreibe Tests (Golden Dataset für LLM-Teil, Unit-Tests für Logik)
    \item Erstelle Dockerfile und docker-compose.yml
    \item Schreibe README mit Architektur-Diagramm und Setup-Anleitung
    \item Teile das Projekt auf GitHub oder in deinem Portfolio
\end{enumerate}



\section{Ressourcen}

\begin{itemize}[leftmargin=*]
    \item \textbf{OpenAI Quickstart}: \href{https://platform.openai.com/docs/quickstart}{platform.openai.com/docs/quickstart}
    \item \textbf{LangChain Tutorials}: \href{https://python.langchain.com/docs/tutorials/}{python.langchain.com/docs/tutorials}
    \item \textbf{Streamlit Docs}: \href{https://docs.streamlit.io}{docs.streamlit.io}
    \item \textbf{DeepEval}: \href{https://github.com/confident-ai/deepeval}{github.com/confident-ai/deepeval}
    \item \textbf{ChromaDB}: \href{https://docs.trychroma.com}{docs.trychroma.com}
    \item \textbf{FastAPI}: \href{https://fastapi.tiangolo.com}{fastapi.tiangolo.com}
\end{itemize}


